
\def\PUDcodigo{ECA233}

\def\PUDcodacad{04507.31}

\def\PUDdisciplina{Controle I}

\def\PUDsemestre{S6}

\def\PUDhoras{80}

%NOVOS ---------------------------------------------------
\def\PUDhorasTeoria{72}
\def\PUDhorasPratica{8}

\def\PUDinterECA{ 
\hyperref[PUD:ECA210]{Circuitos Elétricos I (04507.20)}

\hyperref[PUD:ECA202]{Controle II (04507.40)}

\hyperref[PUD:ECA235]{Controle III (04507.58)}

\hyperref[PUD:EME132]{Inteligência Computacional Aplicada (04507.61)}
}

\def\PUDatuaisECA{---}

\def\PUDrecursos{Projetor multimídia; Quadro branco e pincel; Aulas em laboratório de informática.}

%----------------------------------------------------------

\def\PUDcreditos{4}

\def\PUDprerequisito{---}

\def\PUDpreacad{---}

\def\PUDprerequisito{ \hyperref[PUD:ECA220]{ECA220} }

\def\PUDpreacad{ \hyperref[PUD:ECA215]{Circuios Elétricos II (04507.24)}}

\def\PUDobjetivos{Familiarizar-se com os principais conceitos na área de controle de sistemas lineares no tempo contínuo, em especial o controle clássico, mas com uma introdução ao controle moderno. Compreender os conceitos básicos dos sistemas de controle em malha fechada e do projeto de controladores pelo método do lugar das raízes. Conhecer ferramentas computacionais concernentes
aos tópicos supracitados. 
%Conhecer os conceitos de variável de controle e processos.  Compreender e analisar os sistemas de controle em malha aberta, com retroação e em malha fechada.
}

\def\PUDementa{
Introdução ao Controle Clássico; Modelagem de Sistemas no Domínio da Frequência; Modelagem no Domínio do Tempo; Resposta e Medidas de Desempenho no Domínio do Tempo; Estabilidade; Erro em Regime Permanente; Controladores Proporcionais, Integrais e Derivativos; Técnica do Lugar Geométrico das Raízes; Projeto de Controladores via Lugar Geométrico das Raízes.
%Sistemas de controle clássicos;  transformadas de Laplace;  Modelagem de Sistemas Físicos;  Sistemas dinâmicos;  Respostas Transitórias;  Lugar das Raízes; Diagramas de Bode;  Espaço de Estados.
}

\def\PUDconteudo{ & Introdução ao Controle Clássico: histórico dos sistemas de controle, conceitos e definições básicas, sistemas de malha aberta, sistemas de malha fechada, exemplos de projetos em sistemas de controle. 
%Introdução ao controle clássico;  
\\ 
 & Modelagem de Sistemas no Domínio da Frequência: função de transferência, representação em diagrama de blocos, funções de transferência em circuitos elétricos, sistemas mecânicos translacionais, sistemas mecânicos rotacionais e sistemas eletromecânicos, circuitos elétricos análogos.
 %Transformada de Laplace;  Expansão em frações parciais;  Solução de equações diferenciais LIT.  transformada inversa de Laplace.  Propriedades.  Análise de sistemas LIT  usando a transformada de Laplace. 
\\ 
 & Modelagem no Domínio do Tempo: representação geral no espaço de estados, exemplos de aplicação, conversão de uma função de transferência para o espaço de estados, conversão do espaço de estados para uma função de transferência.
 %Modelagem de Sistemas Físicos: Sistema de Aquecimento;  Sistema de Nível de Líquido;  Sistema de Vazão de Líquido;  Sistema de Eletro-mecânico;  Estimação de Sistemas (Mínimos Quadrados); 
\\ 
 & Resposta no Domínio do Tempo: polos, zeros e resposta completa de um sistema, sistemas de primeira ordem, sistemas de segunda ordem, sistemas de ordem superior, medidas de desempenho em sistemas de segunda ordem,  solução via transformada de Laplace de equações de estado, solução no domínio do tempo de equações de estado.
 %Sistemas dinâmicos.  Função de transferência e resposta ao impulso;  Diagrama de blocos.  Resposta transitória de sistemas de primeira e seginda ordem.  Sistemas de ordem superior.  critério de estabilidade. 
\\ 
 & Estabilidade e Erro em Regime Permanente: introdução, critério de Routh-Hurwitz, estabilidade no espaço de estados, erro em regime permanente, classificação dos sistemas de controle quanto ao erro, constantes de erro de posição, de velocidade e de aceleração.
 %Método do lugar das raízes.  Construção de diagramas LR.  Lugar das raizes para sistemas com atraso de transporte;  Diagramas de contorno;  
\\ 
 & Controladores: ações de controle proporcional (P), integral (I), derivativo (D), proporcional-integral (PI), porporcional-derivativo (PD) e proporcional-integral-derivativo (PID), exemplos de aplicação.
 %Análise no domínio da frequência.  Diagramas de Bode;  gráficos polares;  Critérios de estabilidade de Nyquist.  Resposta de frequência em malha fechada. 
\\ 
 & Lugar Geométrico das Raízes: introdução, definição e propriedades, esboço do lugar geométrico das raízes (LGR), refinamento do LGR. 
 %Controladores PID.  Sintonia de controladores PID.  
\\ 
 & Projeto Via Lugar Geométrico das Raízes: melhorando o erro em regime permanente via compensação em cascata, melhorando a resposta transitória via compensação em cascata, melhorando o erro em regime permanente e a resposta transitória, exemplos de projeto.
 %Sistemas de controle no espaço de estados.  Representação de funções.  Soluções de equações de Estado e análise matricial. Controlabilidade e Observabilidade. 
 }

\def\PUDmetodo{
Aulas teóricas expositivas com auxílio de recursos audio-visuais, aulas práticas com simulações computacionais e com protótipos reais.
%Aulas expositórias que podem ser teóricas e/ou práticas, onde as práticas no laboratório serão marcadas no decorrer da disciplina.
}

\def\PUDavalia{
As avaliações de conhecimento ocorrerão no decorrer da disciplina na forma de provas teóricas, como também de trabalhos/projetos a serem apresentados na forma de relatórios e/ou seminários.
%A avaliação será feita com aplicação de provas teóricas e/ou práticas, além da possibilidade de inclusão de trabalhos e seminários no decorrer da disciplina.
}

\def\PUDbibbasica{ & \href{www.biblioteca.ifce.edu.br/index.asp?codigo_sophia=63228}{Ogata}, Katsuhiko.  Engenharia de controle moderno.  5º ed.  São Paulo, SP: Pearson Prentice Hall, 2010.  809p.  ISBN: 9788576058106. 
\\ 
 &  \href{www.biblioteca.ifce.edu.br/index.asp?codigo_sophia=62576}{Maya}, Paulo A.; Leonardi, Fabrizio.  Controle essencial.  2º ed.  São Paulo, SP: Pearson Education do Brasil, 2014.  347p.  ISBN: 9788543002415.
\\ 
 & \href{www.biblioteca.ifce.edu.br/index.asp?codigo_sophia=63123}{Nise}, Norman S.  Engenharia de sistemas de Controle.  6º ed.  Rio de Janeiro, RJ: LTC, 2012.  745p.  ISBN: 9788521621355. }

\def\PUDbibcomplementar{ &  \href{www.biblioteca.ifce.edu.br/index.asp?codigo_sophia=62612}{Franklin}, Gene F.; Powell, J. David.  Sistemas de Controle para Engenharia.  6º ed.  Porto Alegre, RS: Bookman, 2013.  702p.  ISBN: 9788582600672.
%&  DORF, Richard C.;  BISHOP, Robert H. , Modern Control Systems.  11 ed.  New Jersey: Pearson Prentice Hall.  2009.  1018p.  ISBN 9780132270281
%&  KOROGUI, Rubens H.;  GEROMEL, José C.  Controle Linear de sistemas Dinâmicos. São Paulo: Ed. Edgard Blucher.  2011.  363p.  ISBN : 9788521205906.
\\ 
 &  \href{www.biblioteca.ifce.edu.br/index.asp?codigo_sophia=24226}{Simões}, Marcelo Godoy.  Controle e modelagem fuzzy.  2º ed.  São Paulo, SP: Blucher : FAPESP, 2007.  ISBN: 9788521204169.
\\
 &  \href{www.biblioteca.ifce.edu.br/index.asp?codigo_sophia=24230}{Campos}, Mario Cesar M. Massa de.  Controles típicos de equipamentos e processos industriais.  São Paulo, SP: Edgard Blücher, 2006.  ISBN: 9788521203988.
\\
 & \href{	biblioteca.ifce.edu.br/index.asp?codigo_sophia=24297}{Alves}, José Luiz Loureiro. Instrumentação, controle e automação de processos. Rio de Janeiro, RJ: LTC, 2005. 270p. ISBN: 9788521614425.
\\
& \href{biblioteca.ifce.edu.br/index.asp?codigo_sophia=24119}{Bolton}, William. Instrumentação e controle: sistemas, transdutores, condicionadores de sinais, unidades de indicação. Volume único. Curitiba, PR: Hemus, 2005. 197p. ISBN: 852890119-X. 
 %\href{www.biblioteca.ifce.edu.br/index.asp?codigo_sophia=56358}{Oppenheim}, Alan V.; Willsky, Alan S; Nawab, Syed Hamid.  Sinais e Sistemas.  E-book.  	Pearson.  594p.  ISBN: 9788576055044.
\\
 &  \href{www.biblioteca.ifce.edu.br/index.asp?codigo_sophia=23877}{Rosário}, João Maurício.  Princípios de mecatrônica.  São Paulo, SP: Prentice Hall, 2005.  356p.  ISBN: 9788576050100. }

\def\PUDautor{Adriano Pereira}

\def\PUDdata{2013-04-27}

\def\PUDrevisao{3}

\def\PUDrevisor{José Daniel de Alencar Santos}

\def\PUDdatarevisao{2019-05-15}

\PUDcorpo
